\documentclass{book}
\usepackage{listings}
\usepackage{hyperref}
\usepackage{xcolor}
\usepackage{geometry}
\geometry{a4paper, margin=1in}
\usepackage{parskip}

\lstset{
    basicstyle=\ttfamily\small,
    backgroundcolor=\color{black!5},
    frame=single,
    breaklines=true,
    keywordstyle=\color{blue}\bfseries,
    commentstyle=\color{green!50!black},
    stringstyle=\color{red},
    showstringspaces=false
}

\title{Frust Language Reference}
\author{LagDaemon Research}
\date{\today}

\begin{document}
\maketitle
\tableofcontents

\chapter{AGENTS.md}

Rules for any AI agent (Claude, Codex, Gemini, etc.) working in this repo.

This repo was migrated out of a larger personal R&D monorepo
(\texttt{lagdaemon-tech-research}) on 2026-08-24 once Frust outgrew being one
experiment among many - full commit history preserved via
\texttt{git filter-repo}. Everything below carries forward from that repo's
own hard-won \texttt{AGENTS.md}, adjusted for this repo now being Frust-only.

\section{Platform target - read this first}

\textbf{Windows only. No Linux/Unix/cross-platform consideration, ever, for
any part of this project — the language, the compiler, the IDE, the
standard library — unless the user explicitly asks for it in that
specific moment.} Do not pause, defer, or scope a feature down
because "it wouldn't be portable" or "there's no platform-conditional
compilation mechanism yet." Do not add \texttt{#ifdef}/\texttt{#cfg}-style
portability branches, do not write code with an eye toward a future
Linux port, do not note "Windows-only for now" as a caveat implying
it should change later. Build it for Windows and ship it.

The user has stated this directly, repeatedly, across sessions in the
old repo ("I have said this umpteen times... it does not need to be
compatible with unix, stop that, its not a consideration for you now"
\begin{itemize}
\item 2026-08-22). Treat it as settled, not open for reconsideration. If
\end{itemize}
you notice yourself about to hedge on cross-platform grounds, that is
the signal to stop and just build the Windows version.

\section{Build}

\begin{itemize}
\item \textbf{Debug builds, not Release.} Build/run with \texttt{--config Debug}, not \texttt{Release}, unless explicitly told otherwise for a specific test.
\item \textbf{Single-core builds only.} Never pass \texttt{/m} or \texttt{/maxcpucount} to MSBuild (or equivalent parallel-build flags to other build tools) on this machine — this is the user's own machine and they need it usable while a build runs. Plain \texttt{MSBuild.exe solution.sln /t:target /p:Configuration=Debug}, no parallelism flag, every time.
\item Frust core configures via a CMake preset (\texttt{windows-vcpkg}) that supplies the LLVM install root plus WinFlexBison/JUCE paths for this machine - \texttt{cmake --preset windows-vcpkg} from \texttt{projects/01_language_paradigms/02_functional}, not a bare \texttt{cmake ..}.
\item Never configure Frust with vcpkg's CMake toolchain file and never run \texttt{vcpkg install} from this repo as part of a normal Frust build. The Windows preset consumes the already-built LLVM install at \texttt{D:/000 Creation Suite/apps/CreationEngine/vcpkg_installed/x64-windows}; \texttt{CMakeLists.txt} must fail if LLVM is not found inside that existing install tree.
\item Other subprojects (\texttt{09_frust_plugin_host}, \texttt{10_node_compiler}, \texttt{02_juce_language_host}) pull in Frust core via \texttt{add_subdirectory} using a relative path up to \texttt{01_language_paradigms/02_functional} - the \texttt{projects/} layout must stay intact for these to resolve.
\end{itemize}

\section{Project naming}

\begin{itemize}
\item New project folders under \texttt{projects/} use kebab-case with \texttt{-} separators, not spaces or numeric prefixes. Example: \texttt{frust-linalg}, not \texttt{13_frust_linalg} or \texttt{Frust Linalg}.
\end{itemize}

\section{Standard workflow (every change, no exceptions)}

1. Make the change.
2. Build it (Debug config) and actually run/verify it — don't claim done on compile-success alone. For a language/codegen change, write a real test with a hand-predicted expected result BEFORE running it, then run and confirm the match.
3. Run the full regression sweep (\texttt{09_frust_plugin_host}'s example suite, \texttt{02_juce_language_host} rebuild) when the change touches core codegen/grammar paths - not just the one new test.
4. \texttt{git status} / \texttt{git diff} — check what's actually changed before staging.
5. Commit, with a real message explaining why, not just what.
6. Push to \texttt{origin} — the user has standing authorization to push routinely ("push early, push often"); still use judgment on force-push/branch-deletion style destructive ops, those always need explicit sign-off.
7. If it's a submodule (\texttt{projects/06_frust_library}, its own repo at \texttt{wwestlake/frust-library}), commit+push inside the submodule first, THEN commit+push the parent repo's pointer bump. Two separate pushes, always in that order.
8. If you told the user a rule/process applies going forward, write it into this file, not just into your own private memory. This file is what's authoritative and visible in the repo — private memory is a supplement, never a substitute.
9. Keep \texttt{LANGUAGE_GAPS.md} (\texttt{projects/01_language_paradigms/02_functional/}) current the moment an item's status changes - this document existing isn't enough on its own, the standing failure mode it was built to fix is real progress happening underneath it while the doc goes stale.
10. Wiki/tutorial/doc examples are a real use-case test of the compiler, not just prose: run every example through \texttt{frust_compiler.exe} (hand-predicted output vs. actual) before it goes in a page. If an example reveals a real bug - not just an already-documented gap - stop the doc work immediately, log it in \texttt{LANGUAGE_GAPS.md}'s "CRITICAL BUGS" section if it's silent/incorrect-output-class (jumps the queue ahead of normal numbered gaps), fix it, verify the fix, then resume using the corrected pattern. Don't quietly route around a found bug with a "safer" example instead.

\section{Git discipline}

\begin{itemize}
\item \textbf{Commit AND push proactively, both, every time.} After every verified build/change: commit right then, push to \texttt{origin}/\texttt{main} right after. Don't wait to be asked for either one.
\item \textbf{This repo has no deployment - it's just a git repo.} \texttt{main} is the working branch; pushing straight to it carries no CI/CD or production risk. Treat pushing to \texttt{main} as routine, not something to hesitate over.
\item Never force-push, never \texttt{--no-verify}, never skip hooks without being told to.
\item Never summarize a set of tracked work (a gap list, a checklist) as "done"/"closed"/"complete" if even one item in it is genuinely PARTIAL - the summary sentence must reflect the worst status in the set, not the majority. Say "N of M done, rest open" plainly.
\end{itemize}

\section{Communication}

\begin{itemize}
\item \textbf{Don't silently downgrade a possibly-transient failure into "it's gone" before an expensive fallback.} Say the assumption out loud first, especially before long rebuilds.
\item \textbf{State the target machine/file directly instead of asking the user to disambiguate} when context already makes it obvious.
\item \textbf{Don't ask the user to make implementation-level technical decisions} — make the call and proceed. Only surface real goals/priorities as questions.
\item \textbf{Discuss real language/tooling design choices before implementing them} — don't just build and present as done. Reserve this for genuine architectural calls (a new type-system feature, a new codegen representation), not mechanical fixes.
\item \textbf{Critical handoff info (passwords, IPs, key facts) gets its own short, clearly labeled block} — never buried in a paragraph.
\item \textbf{Never trigger UAC/sudo/admin elevation beyond already-agreed scope without discussing first}, every time, no exceptions.
\end{itemize}

\section{Toolchain}

\begin{itemize}
\item LLVM 18.1.6 via vcpkg, WinFlexBison, and JUCE are all machine-level installs referenced by absolute path in \texttt{CMakePresets.json} - not part of this repo, don't try to vendor them in.
\end{itemize}


\chapter{Handoff to Codex — Claude session, 2026-08-24 through 2026-08-25}

Date: 2026-08-25
From: Claude, repo \texttt{wwestlake/FrustLang}, branch \texttt{master}

This is a pre-emptive handoff: Claude is nearing a weekly token limit and may
go offline mid-task. Per a hard-learned lesson from a previous session (a
full day's work sat unpushed locally and had to be recovered by hand), the
standing rule now is \textbf{push everything the moment it's real, don't wait for
a natural stopping point.} Everything described below as "done" is genuinely
committed and pushed to GitHub - verify with \texttt{git log}/\texttt{gh pr list} if in
doubt, don't take this file's word for it blindly.

\section{Open PRs / issues (nothing urgent, but worth knowing about)}

\begin{itemize}
\item \textbf{frust-library#1} (\texttt{fix/wire-missing-core-modules}) and \textbf{frust-library#2}
\end{itemize}
  (\texttt{feature/jit-eval-wrapper}, based on top of #1) - both open, both pushed,
  both real fixes to the \texttt{06_frust_library} submodule (core pod wiring,
  2 real pre-existing crash bugs, and a new \texttt{jit_eval_f32} wrapper). #2
  includes #1's commits, so merging #2 alone (into frust-library's \texttt{main})
  effectively closes both.
\begin{itemize}
\item \textbf{FrustLang#9} (issue, not PR) - a verification checklist for the big
\end{itemize}
  PR #8 that just merged (see below). Read this before assuming anything
  in that PR's scope is fully proven.
\begin{itemize}
\item FrustLang PRs #1-#8 are all merged into \texttt{master} already - nothing else
\end{itemize}
  open on this repo as of this writing.

\section{What just happened (PR #8, merged into master)}

A large, connected chain of real fixes, all triggered by trying to write an
honest GitHub wiki tutorial "First Program" lesson using real Frust (no raw
\texttt{extern fn}) instead of hello-world boilerplate. Full detail is in PR #8's
own description and its commit message - short version:

1. \textbf{New \texttt{import <pod>, "<version>";} keyword} - real cross-pod source-level
   imports (\texttt{frust.y}/\texttt{frust.l}/\texttt{AST.h}), distinct from the pre-existing
   \texttt{use somepod;} (which parses but has always been a complete no-op).
   \texttt{"current"} reads the version already declared in the importing pod's own
   \texttt{frate.json}. Resolution logic lives in \texttt{frate}'s own file-collection code
   (\texttt{collectImportedPodFiles}, \texttt{projects/05_frate/src/main.cpp}), mirroring
   how \texttt{use self::X;} resolution already worked there.
2. \textbf{Extracted \texttt{frust_runtime}} - \texttt{frust_print_str}/\texttt{format_*}/\texttt{buf_*}/
   \texttt{ast_*}/\texttt{jit_eval_f32} used to be duplicated verbatim in both
   \texttt{frust_compiler}'s and the IDE's \texttt{Main.cpp}. Now one real static library
   (\texttt{projects/01_language_paradigms/02_functional/Runtime.cpp} + its
   CMakeLists.txt target) both link.
3. \textbf{Fixed \texttt{frate}'s AOT link step}, which never actually produced a
   working standalone executable before this. Removed \texttt{/ENTRY:main} (was
   bypassing CRT startup entirely - see PR #8's commit message for the full
   "how other languages do this" reasoning), wired in \texttt{frust_runtime.lib}
   and (conditionally) \texttt{frust_plugin_host.lib} + its full LLVM/Windows
   dependency chain, switched to a linker response file (the raw argument
   list blew past \texttt{cmd.exe}'s line-length limit).
4. \textbf{\texttt{frate cache-dir} command} + configurable pod cache location
   (\texttt{projects/05_frate/include/frate/FrateCache.h}/\texttt{.cpp}) - was hardcoded
   to \texttt{%APPDATA%} (always writes to C:), now resolves env var → persisted
   \texttt{frate_settings.json} next to the exe → default next to the exe.
5. \textbf{\texttt{examples/tutorial-first-program/}} - the wiki's real "First Program"
   lesson kept in-repo, wired into \texttt{.github/workflows/release.yml} as a real
   release gate.

\section{What's verified vs. not (see issue #9 for the live checklist)}

\textbf{UPDATE, same session, after the full rebuild finished}: all three
locally-checkable items below are now done, against the exact merged
commit, not a scratch copy:
\begin{itemize}
\item ✅ Full \texttt{frust_plugin_host} regression suite - all 16 examples,
\end{itemize}
  \texttt{ALL_CHECKS_PASSED} on every one.
\begin{itemize}
\item ✅ Fresh IDE Debug build + launch smoke test - builds clean, launches,
\end{itemize}
  stays running.
\begin{itemize}
\item ✅ The literal committed \texttt{examples/tutorial-first-program/} - \texttt{frate run}
\end{itemize}
  against it (not a scratch copy) produces exactly the expected output
  (\texttt{Entrance Hall} / \texttt{A cold stone hall...}), exit 0.

\textbf{Still genuinely open} - the one item that needs a real GitHub Actions
runner, can't be checked locally: a real CI run of
\texttt{.github/workflows/release.yml}'s tutorial-gate step. That only fires on a
\texttt{v*} tag push, which also publishes a real GitHub Release as a side effect
\begin{itemize}
\item not something to trigger casually just to test the step. Whoever pushes
\end{itemize}
the first real \texttt{v0.5.0} tag is the first real end-to-end test of it.

Original verification note (kept for the record): a scratch copy of the
tutorial example was checked first, before the full rebuild finished, and
JIT mode (Hello World, live metaprogramming) was re-confirmed working
right after the \texttt{frust_runtime} extraction, before this PR was pushed.

\textbf{If you're picking this up because Claude ran out of tokens mid-verification:}
check whether a Claude session already ran/reported on the above (look for
recent commits, or a comment/note in issue #9) before redoing it from
scratch. If nothing's there, start with the regression suite - it's the
highest-value check (catches any real regression across the whole plugin
host stack, not just the one new example).

\section{Real gotchas for this repo (save yourself the rediscovery time)}

\begin{itemize}
\item \textbf{Single-core builds only.} Never pass \texttt{/m} or \texttt{/maxcpucount:N>1} to
\end{itemize}
  MSBuild on this machine - has caused real problems before. Always
  \texttt{/maxcpucount:1}.
\begin{itemize}
\item \textbf{MSBuild.exe lives on D:, not the usual C: path}:
\end{itemize}
  \texttt{D:\Program Files\Microsoft Visual Studio\2022\Community\MSBuild\Current\Bin\MSBuild.exe}.
  Same for \texttt{vswhere.exe}: \texttt{C:\Program Files (x86)\Microsoft Visual
  Studio\Installer\vswhere.exe} (that one IS on C:, just MSBuild itself isn't).
\begin{itemize}
\item \textbf{LLVM is NOT installed via a normal vcpkg install on this machine} - it's
\end{itemize}
  at \texttt{D:\000 Creation Suite\apps\CreationEngine\vcpkg_installed\x64-windows},
  a real, trimmed (~11GB) build, also mirrored on S3
  (\texttt{s3://creation-suite-build-cache/vcpkg-cache/llvm-vcpkg-x64-windows.tar})
  for CI. Don't let anything trigger a fresh LLVM build from source - it's
  slow and unnecessary, the trimmed build already covers everything needed.
\begin{itemize}
\item \textbf{Windows only, no Linux/Unix consideration anywhere in this project,
\end{itemize}
  ever}, unless explicitly asked in the moment - a standing, deliberate,
  repeatedly-stated directive.
\begin{itemize}
\item \textbf{\texttt{AGENTS.md}} (repo root) has the full standing rules list - read it
\end{itemize}
  before doing anything nontrivial. Rule 10 in particular: any tutorial/doc
  example is a real use-case test of the compiler - run it through
  \texttt{frust_compiler.exe}/\texttt{frate} for real before publishing it, and if it
  reveals a real bug, stop, fix the bug, then resume - don't quietly route
  around it with a "safer" example.
\begin{itemize}
\item \textbf{\texttt{LANGUAGE_GAPS.md}} (\texttt{projects/01_language_paradigms/02_functional/})
\end{itemize}
  is the authoritative, actively-maintained record of what's actually done
  vs. partial vs. open in the language itself - keep it current the moment
  something's status changes, don't let it go stale.
\begin{itemize}
\item \textbf{\texttt{06_frust_library} is a git submodule}, its own separate repo
\end{itemize}
  (\texttt{wwestlake/frust-library}). Commit/push order matters: submodule first,
  then the parent repo's pointer bump, always in that order, two separate
  pushes.
\begin{itemize}
\item \textbf{PR-based workflow, not direct push} - branch, commit, push the branch,
\end{itemize}
  \texttt{gh pr create}. The user has owner-level bypass on branch protection and
  may merge directly themselves sometimes; don't assume that's the default
  for anyone else.
\begin{itemize}
\item \textbf{The GitHub wiki} (github.com/wwestlake/FrustLang/wiki) has a partial
\end{itemize}
  tutorial published (Getting Started, Building from Source) that's
  mid-rewrite as a build-along "single-user dungeon crawl" project (design
  captured in a Claude plan file, not in this repo) - the
  \texttt{examples/tutorial-first-program/} pod is that rewrite's actual "First
  Program" lesson, kept in sync with whatever the wiki page says. If you
  touch the wiki, keep both in sync.

\section{Where things are headed (not started, just named so it isn't a surprise)}

\begin{itemize}
\item The dungeon-crawl tutorial rewrite is far from finished - only "Getting
\end{itemize}
  Started"/"Building from Source"/"First Program" exist so far, of a
  planned ~13-lesson arc (control flow, interfaces, generics, \texttt{shared}
  ownership, save/load, an \texttt{extern}-FFI lesson calling a small planned C++
  terminal-UI library named "Agent 86", the IDE, plugins-as-mods, and an
  optional live-metaprogramming capstone).
\begin{itemize}
\item A \texttt{v0.5.0} GitHub Release hasn't actually been tagged/published yet - the
\end{itemize}
  CI pipeline (\texttt{.github/workflows/release.yml}) that would build and publish
  it exists and is believed correct, but has never been triggered by a real
  tag push.
\begin{itemize}
\item \textbf{\texttt{v0.5.0} tag pushed} (2026-08-25, same session as everything above) -
\end{itemize}
  the very first real run of the CI release pipeline, triggered live. Check
  https://github.com/wwestlake/FrustLang/actions for the result if it isn't
  already reflected in a comment on issue #9 or a later handoff note - this
  session's token budget ran out before it could see the run through to
  completion (a real 20-60+ minute Windows CI build, first time ever run).
  The workflow itself doesn't depend on this session staying alive - it's
  running entirely on GitHub's own infrastructure regardless.
\begin{itemize}
\item \textbf{New idea, spec being developed in a separate session (not this repo
\end{itemize}
  yet)}: a .NET host for Frust, so it could be embedded in a .NET app.
  Assessed as genuinely feasible - \texttt{frust_plugin_host}'s API is already a
  clean \texttt{extern "C"} C ABI, exactly what .NET's P/Invoke needs. The one real
  gap: it's only ever been built as a static \texttt{.lib} so far, never a \texttt{.dll} -
  P/Invoke needs a real dynamic library to load at runtime, so step one
  would be adding a DLL build target (a modest CMake addition, not a
  redesign, since the API surface is already right). Not started - whoever
  picks this up should look for the actual spec first (being written
  elsewhere) before designing anything.


\chapter{FrustLang}

Frust is a statically-typed, compiled language built on LLVM (JIT and
AOT), designed around zero-overhead control: real generics
(monomorphized, not type-erased), real algebraic data types with
pattern matching, \texttt{shared} reference counting with automatic drop, real
closures, real interface dispatch - all compiled, never interpreted.

This repo is the whole Frust ecosystem: the compiler, the standard
library, a plugin/host system for embedding Frust in other
applications, a JUCE-based IDE, and supporting tooling (VS Code
extension, language server, a graph-to-source node compiler). It was
split out of a larger personal R&D monorepo once Frust outgrew being
one experiment among many - full commit history was preserved across
the move.

\section{What's actually working}

Tracked honestly in
\texttt{LANGUAGE_GAPS.md} -
every feature below moved from "designed" to "implemented" to
"verified with a real hand-predicted test" before being called done,
and the gaps document says exactly what's still open.

\begin{itemize}
\item \textbf{Real generics, monomorphized}: \texttt{struct Box<T>} and generic free
\end{itemize}
  functions (\texttt{fn identity<T>(x: T) -> T}, called via explicit
  \texttt{identity::<i64>(5)} turbofish syntax) each get their own concrete
  LLVM type/function per instantiation - no boxing, no vtables unless
  you asked for one.
\begin{itemize}
\item \textbf{Real algebraic data types}: \texttt{enum} is a genuine discriminated
\end{itemize}
  union (F#/Rust-style, not a C-style tag-only enum) - variants carry
  typed payloads, including nested enums/structs and generics
  (\texttt{enum Choice<A,B> { Left(A), Right(B) }}). \texttt{match} does real
  recursive pattern matching, reaching through several variant layers
  in one arm, with required exhaustiveness (a compile error if a
  variant - or a \texttt{_} wildcard - is missing).
\begin{itemize}
\item \textbf{\texttt{Result<T,E>} / \texttt{Option<T>}} (\texttt{enum Result<T,E> { Ok(T), Err(E) }})
\end{itemize}
  built as ordinary Frust code on top of \texttt{enum}, not a compiler
  intrinsic - \texttt{match} enforces both cases are handled, no sentinel
  field to misread.
\begin{itemize}
\item \textbf{Real interface dispatch} (\texttt{interface} / \texttt{impl X for Y}) via fat
\end{itemize}
  pointers, no hidden allocation.
\begin{itemize}
\item \textbf{\texttt{shared}} with real strong reference counting and automatic
\end{itemize}
  scope-exit drop - deliberately designed to fail toward *leaking*,
  never toward a double-free, when ownership can't be proven
  statically. \texttt{own}/\texttt{raw} pointers exist alongside it for explicit,
  zero-overhead control.
\begin{itemize}
\item \textbf{Closures} as \texttt{{ code, env }} fat pointers, capture-by-value,
\end{itemize}
  fully compiled.
\begin{itemize}
\item \textbf{Live metaprogramming} - \texttt{quote}/\texttt{unquote}/\texttt{build_time} blocks
\end{itemize}
  that run and splice real AST at build time.
\begin{itemize}
\item Growable collections (\texttt{Vector<T>}), real pointer arithmetic/raw
\end{itemize}
  dereference, multi-file plugin support, and more.

\section{Repo layout}

\begin{lstlisting}
projects/
  01_language_paradigms/02_functional/   Frust core: grammar, AST, Codegen, the compiler
  02_juce_language_host/                 JUCE-based IDE, built on frust_plugin_host
  03_creation_dock/                      JUCE docking UI (IDE dependency)
  04_ai_provider/                        AI-provider abstraction (IDE dependency)
  05_frate/                              Frust's package manager
  06_frust_library/                      Standard library (its own submodule/repo: frust-library)
  07_vscode_frust/                       VS Code extension
  08_frust_lsp/                          Language server
  09_frust_plugin_host/                  Hot-reloadable plugin/host system - the extension
                                          mechanism the IDE (and any other app) uses to run
                                          JIT-compiled Frust plugins
  10_node_compiler/                      Graph JSON -> real Frust source; the compiler half of
                                          a planned Blueprint-style visual node layer
                                          (see NODE_LANGUAGE_DESIGN.md)
\end{lstlisting}

Each subproject with its own README has more detail - see
\texttt{09_frust_plugin_host/README.md}
and
\texttt{10_node_compiler/README.md}
in particular.

\section{Building}

Windows only, by deliberate standing decision - no cross-platform
consideration anywhere in this codebase.

Frust core builds via a CMake preset that supplies the LLVM/WinFlexBison/JUCE
paths for this machine's toolchain. On Windows this deliberately points at an
already-installed \texttt{vcpkg_installed/x64-windows} tree and rejects the vcpkg
CMake toolchain file, so configuring or building Frust cannot kick off a
fresh LLVM build.

\begin{lstlisting}
cd projects/01_language_paradigms/02_functional
cmake --preset windows-vcpkg
cmake --build build --config Debug
\end{lstlisting}

Other subprojects (\texttt{09_frust_plugin_host}, \texttt{10_node_compiler},
\texttt{02_juce_language_host}) each pull in Frust core via \texttt{add_subdirectory}
using a relative path (\texttt{../01_language_paradigms/02_functional}) - the
\texttt{projects/} layout above must stay intact for those relative paths to
resolve.

\texttt{06_frust_library} is a git submodule - clone with
\texttt{git clone --recurse-submodules}, or run
\texttt{git submodule update --init} after a plain clone.

\section{Design philosophy}

"Zero-overhead, aiming for the iron" - stated directly in
\texttt{FRUST_LANG_SPEC.md}
and treated as a real constraint, not a slogan: it's the reason
generics are monomorphized rather than type-erased, and the reason
\texttt{shared}'s reference counting has a hard scope limit (own's
automatic-free was deliberately left unimplemented rather than risk a
double-free in code that works today) instead of reaching for a
general-purpose GC.


\chapter{Frust — Reddit/Quora post tracker}

Tracks what's been written and posted for Frust on r/FrustLang and Quora, so we're not
duplicating topics or losing track of what's already out there. Add a new entry below
every time a post is drafted or actually published.

\section{Posted}

\subsection{r/FrustLang — general Frust intro}
\begin{itemize}
\item \textbf{Status:} posted
\item \textbf{Platform:} Reddit, r/FrustLang
\item \textbf{Topic:} general introduction to Frust — what it is, what's actually working.
\item \textbf{Note:} drafted and posted earlier in the session that did the repo migration; exact
\end{itemize}
  text wasn't saved to a file at the time, so it isn't reproduced here verbatim.

\subsection{r/FrustLang — why \texttt{shared<T>} works the way it does}
\begin{itemize}
\item \textbf{Status:} posted
\item \textbf{Platform:} Reddit, r/FrustLang
\item \textbf{Topic:} design-philosophy piece on \texttt{shared<T>} reference counting — deterministic
\end{itemize}
  drop, no GC, why it was built this way. Written using "Frust does X" phrasing
  throughout (not "my language"), per an explicit correction during drafting.
\begin{itemize}
\item \textbf{Note:} exact text wasn't saved to a file at the time, so it isn't reproduced here
\end{itemize}
  verbatim.

\subsection{Quora (Creative Programming space) — "How does Frust handle pointers?"}
\begin{itemize}
\item \textbf{Status:} posted
\item \textbf{Platform:} Quora, Creative Programming space (117.4K followers)
\item \textbf{Topic:} answers "How does Frust handle pointers?" with both real Frust usage code
\end{itemize}
  and real internal \texttt{Codegen.h} implementation snippets (\texttt{sharedHeaderPtr},
  \texttt{dropSharedLocal}, \texttt{retainSharedLocal}). Links to the FrustLang repo and
  https://www.reddit.com/r/FrustLang/.
\begin{itemize}
\item \textbf{Note:} exact text wasn't saved to a file at the time, so it isn't reproduced here
\end{itemize}
  verbatim.

\section{Drafted, not yet posted}

\subsection{5-post series — main serious features (2026-08-24)}
\begin{itemize}
\item \textbf{Status:} drafted, not yet posted
\item \textbf{Platform:} intended for r/FrustLang (5 separate short posts)
\item \textbf{Topic:} one language feature per post, ~2 paragraphs + a short verified code
\end{itemize}
  example each. Syntax in every example was pulled directly from the grammar
  (\texttt{frust.y}) and \texttt{Codegen.h}, not reconstructed from memory.

---

#### 1. Real interface dispatch — not duck typing, actual dynamic dispatch

Frust has \texttt{interface}/\texttt{impl X for Y}, and it's not sugar — a function that takes an
interface-typed parameter genuinely dispatches through a vtable at runtime, so the same
call site can run different code depending on what concrete type got passed in. No
inheritance hierarchy required, no "does this struct happen to have the right method
names" duck typing — you declare conformance explicitly with \texttt{impl}, and the compiler
enforces it.

\begin{lstlisting}
interface Automation {
    fn tick(input: f32) -> f32;
}

impl Automation for Motor {
    fn tick(input: f32) -> f32 = { input * 2.0 }
}

fn run_step(a: Automation, x: f32) -> f32 = { a.tick(x) }
\end{lstlisting}

---

#### 2. \texttt{shared} — reference counting without a garbage collector

Frust's memory model is explicit, not GC'd: \texttt{own} is single-ownership with automatic
drop at scope end, and \texttt{shared} is deterministic reference counting — retain on
copy, release on scope exit, deallocate the instant the count hits zero. No
stop-the-world pause, no background collector thread, and no ambiguity about when a
destructor runs. You always know exactly when cleanup happens by reading the code.

\begin{lstlisting}
fn main() -> i64 = {
    let a: shared Counter = shared Counter { value: 0 };
    let b = a;                 // retain - refcount 2
    // both drop at end of scope, refcount hits 0, real free happens here
    0
}
\end{lstlisting}

(Corrected 2026-08-24: the type annotation is \texttt{shared Counter}, a bare
prefix keyword before the plain type name — \texttt{shared<Counter>} isn't
valid Frust syntax. Caught before this post went out; see
\texttt{LANGUAGE_GAPS.md}'s "CRITICAL BUGS" section for the full story of how
this was found.)

---

#### 3. Generics — real monomorphization, not type erasure

\texttt{struct Box<T>} and \texttt{fn identity<T>(x: T) -> T} are genuine generics: each concrete
instantiation gets its own compiled, specialized version (same strategy Rust and C++
templates use), not a boxed/erased runtime representation. Turbofish syntax picks the
concrete type explicitly at the call site when it can't be inferred from arguments
alone.

\begin{lstlisting}
fn identity<T>(x: T) -> T = { x }

fn main() -> i64 = {
    let a = identity::<i64>(42);
    let b = identity::<f32>(3.14);
    0
}
\end{lstlisting}

---

#### 4. Live \texttt{quote}/\texttt{unquote}/\texttt{build_time} — metaprogramming that runs while your program is already running

This is the one people don't expect: \texttt{build_time { }} blocks can \texttt{quote { }} real AST
nodes (with \texttt{unquote(expr)} splicing live runtime values in as literals), hand that
AST to the compiler's own JIT, and get back a freshly-compiled function pointer — all
from inside a program that's already executing. It works because a running \texttt{.frust}
program shares a process with the same compiler pipeline that built it, so "generate
code from live data and compile it in on the spot" is just calling back into that
pipeline, not a separate execution model bolted on.

\begin{lstlisting}
fn build_line_eval(m: f32, b: f32) -> ASTExpr = build_time {
    quote {
        (unquote(m) * x) + unquote(b)
    }
}

extern fn frust_jit_eval_f32(ast: ASTExpr, x: f32) -> f32;

fn main() -> i64 = {
    let ast = build_line_eval(2.0, 1.0);   // y = 2x + 1, specialized just now
    let y = frust_jit_eval_f32(ast, 5.0);  // JIT-compiles and runs it: 11.0
    0
}
\end{lstlisting}

---

#### 5. Embeddable by design — Frust runs inside your app, not just next to it

Frust compiles to native x86, but it also JITs and links in-process, which means it's
built to be *embedded*: the JUCE-based IDE in this repo hosts Frust plugins directly
in its own process, calling into JIT-compiled Frust functions like any other C
function pointer. That makes Frust usable as a real scripting/extension layer for a
host application — not a separate process you shell out to, not a subprocess you pipe
JSON through, but code that runs in the same address space as your app.

\begin{lstlisting}
// host app (C++), loading and calling a Frust plugin in-process
FrustPluginHandle plugin = frust_plugin_load("linter.frust");
void* fn = frust_plugin_get_fn(plugin, "check_line");
auto check = reinterpret_cast<int64_t(*)(const char*)>(fn);
int64_t issues = check("// TODO: fix this");
\end{lstlisting}



\end{document}
